\section{The switched rows $S_{11},S_{12},S_{13}$}

This section proves a joint lower bound for the switched rows.

\subsection{The source and its finite signs}

This subsection describes the source and proves the finite-source sign lemma.
Fix the parameters
\[
 \alpha={100\over1149},\quad\beta={50\over309},\quad
 \lambda={93\over400},\quad\mu={25\over128},\quad u={1\over500},
 \qquad \varepsilon_A=\eta_L=10^{-4}.
 \tag{4.1}
\]
Let $0<\sigma\le10^{-10}$, and put
\[
 \tau_g=\tau=\sigma^9,\quad \rho=1+\tau,\quad L=\log N,\quad
 U_C={C(N)N\over L^2},\qquad \mathcal W=4(1-\varepsilon_A)U_C.
 \tag{4.2}
\]
The D-unit denominator is $\mathcal W$.
All sieve primes belong to $\mathcal P(N)=\{p>2:p\nmid N\}$.
Source primes dividing $N$ enter the final error table,
as do repeated source primes in squarefree identities.
Write
\[
 \mathcal A_{\varepsilon_A}
 =\{e=N-p:p\text{ prime},\ p<(1-\varepsilon_A)N\}.
 \tag{4.3}
\]
Let $t_i=\log_Np_i$.
The source forms, each with $\varepsilon_A N<e<N$, are as follows.
\begingroup
\setlength{\tabcolsep}{3pt}
\begin{longtable}{@{}>{\raggedright\arraybackslash}p{0.24\textwidth}>{\raggedright\arraybackslash}p{0.68\textwidth}@{}}
\toprule
Form & Exact representation and constraints \\
\midrule
$D_{11,3}$ & $e=p_1p_2p_3$, $N^\lambda\le p_1<N^{1/3}$, $p_1<p_2$, $p_1p_2^2<N$, $p_3\ge p_2$ \\
$D_{11,4}$ & $e=p_1p_2p_3p_4$, the same outer conditions, $p_2\le p_3\le p_4$ \\
$D_{12}$ & $e=p_1p_2p_3$, $N^\alpha\le p_1<N^{1/3}\le p_2$, $p_1p_2^2<N$, $p_3\ge p_2$ \\
$D_{13}$ & $e=p_1p_2p_3$, $N^\beta\le p_1<N^\lambda\le p_2$, $p_1p_2^2<N$, $p_1p_2p_3^2\ge N$ \\
\bottomrule
\end{longtable}
\endgroup
On $D_{13}$, the two endpoint orders give distinct representations.
The condition $p_3\ge p_2$ is omitted there.
Let $R_{11}=R_{11,3}+R_{11,4}$, $R_{12}$ and $R_{13}$ count representations, and put
\[
 A_G(e)=(2R_{11}(e)+R_{12}(e)+R_{13}(e))
                  \mathbf1_{\varepsilon_A N<e<N}.
 \tag{4.4}
\]
These forms exhaust the sources.
For $D_{11}$, $5\lambda>1$ allows at most two cofactor primes,
and $\varepsilon_A N>N^{2/3}$ excludes cofactor 1.
A composite $D_{12}$ cofactor gives $e\ge N^{\alpha+1}>N$.
Two $D_{13}$ cofactor factors give $e\ge N$, because each factor is at least $\sqrt{N/(p_1p_2)}$.
The linked source parameters are:
\begin{center}
\begin{tabular}{@{}llll@{}}
\toprule
Form & Weight & HB patterns/factors & Cofactor cutoff \\
\midrule
$D_{11,3}$ & $2$ & $27$ / $6\cdot3$ & $p_2$ \\
$D_{11,4}$ & $2$ & $81$ / $24$ & $p_2$ \\
$D_{12}$ & $1$ & $27$ / $6\cdot3$ & $p_2$ \\
$D_{13}$ & $1$ & $27$ / $6\cdot3$ & $\sqrt{N/(p_1p_2)}$ \\
\bottomrule
\end{tabular}
\end{center}
All four forms use rows G1-0 through G1-3.
Their faces are listed in (4.36) and the boundary-family table.
Their multiplicities satisfy (4.6) and the sharper peel bound below.
Every source coordinate is at least $N^\alpha$.
On $D_{11,3}$ and $D_{11,4}$, each remaining factor is at least
$p_2>p_1\ge N^\lambda$.
On $D_{12}$, $p_3\ge p_2\ge N^{1/3}$.
Only on $D_{13}$ do we use its square-root cutoff. There
$t_2<(1-t_1)/2$ and $t_1<\lambda<1/3$, so
\[
 t_3\ge {1-t_1-t_2\over2}>{1-t_1\over4}>
 {1\over6}>\alpha.
 \tag{4.5}
\]
For absolute estimates,
\[
 A_G(e)\le2(3!+4!)+3!+3!=72<100.
 \tag{4.6}
\]
For the peel, $A_G\le3$.
Ordering gives $R_{11},R_{12}\le1$.
If $R_{11}>0$, all factors are at least $N^\lambda$, so $R_{13}=0$.
Otherwise $D_{13}$ chooses one low-band and one high-band factor,
with at most $lh\le2$ choices, where $l+h=3$ and $l,h\ge1$.
Thus $2R_{11}+R_{12}+R_{13}\le3$, before any union.

Let $I_j(e)$ count the admissible divisor tuples,
retaining $P^-(e/d)\ge z$ and every original coprimality condition.
The identity is
\[
 W_{13}=3I_1+I_2-2I_3-I_4-I_5-I_6-I_7
                 +I_8+I_9+I_{10}-2I_{11}-I_{12}-I_{13}.
 \tag{4.7}
\]
With $\tau_S=57/224$, put
\[
 B(r)=
 \begin{cases}
 \beta,&\beta\le r<\mu,\\
 \min\{\beta,\max\{\alpha,57/64-u-(5/2)r-(11/5)\alpha\}\},
             &\mu\le r<\lambda.
 \end{cases}
 \tag{4.8}
\]
\begingroup
\setlength{\tabcolsep}{3pt}
\begin{longtable}{@{}>{\raggedright\arraybackslash}p{0.1\textwidth}>{\raggedright\arraybackslash}p{0.1\textwidth}>{\raggedright\arraybackslash}p{0.51\textwidth}>{\raggedright\arraybackslash}p{0.21\textwidth}@{}}
\toprule
Row & Weight & Divisor variables and range & Cofactor cutoff \\
\midrule
$S_1$ & $3$ & $d=1$ & $N^\alpha$ \\
$S_2$ & $1$ & $d=1$ & $N^\beta$ \\
$S_3$ & $-2$ & $d=q$, $\alpha\le\log_Nq<\mu$ & $N^\alpha$ \\
$S_4$ & $-1$ & $d=q$, $\mu\le\log_Nq<1/4$ & $N^\alpha$ \\
$S_5$ & $-1$ & $d=q$, $1/4\le\log_Nq<\tau_S$ & $N^\alpha$ \\
$S_6$ & $-1$ & $d=q$, $\tau_S\le\log_Nq<1/3$ & $N^\alpha$ \\
$S_7$ & $-1$ & $d=q$, $\mu\le\log_Nq<\lambda$ & $N^\alpha$ \\
$S_8$ & $1$ & $d=q_1q_2$, $\alpha\le\log_Nq_2<\log_Nq_1<\beta$ & $N^\alpha$ \\
$S_9$ & $1$ & $d=q_1q_2$, $\alpha\le\log_Nq_2<\beta\le\log_Nq_1<\mu$ & $N^\alpha$ \\
$S_{10}$ & $1$ & $d=q_1q_2$, $\mu\le r=\log_Nq_1<\lambda$, $\alpha\le\log_Nq_2<B(r)$ & $N^\alpha$ \\
$S_{11}$ & $-2$ & $p_1,p_2$ in the $D_{11}$ outer domain & $p_2$ \\
$S_{12}$ & $-1$ & $p_1,p_2$ in the $D_{12}$ outer domain & $p_2$ \\
$S_{13}$ & $-1$ & $p_1,p_2$ in the $D_{13}$ outer domain & $\sqrt{N/(p_1p_2)}$ \\
\bottomrule
\end{longtable}
\endgroup

\begin{lemma}[Finite-source signs]
Let $N$ be sufficiently large.
Suppose $e\in(\varepsilon_A N,N)$ is squarefree and coprime to $N$.
If $P^-(e)\ge N^\alpha$ and $\Omega(e)\in\{3,4\}$, then
\[
 W_{13}(e)\le0.
 \tag{4.9}
\]
\end{lemma}
\begin{proof}
The exact sum $\sum t_i$ is kept.
Let $k_i$ count primes in
$[\alpha,\beta)$, $[\beta,\lambda)$, $[\lambda,1/3)$ and $[1/3,1)$.
For a second-band exponent $v$, let $m_B(v)$ count first-band primes below $B(v)$.
The first ten rows give
\[
 B_0=3+\mathbf1_{k_1=0}-2k_1-2k_2-k_3+
                    {k_1\choose2}+\sum_{j=1}^{k_2}m_B(v_j),
 \qquad0\le m_B(v_j)\le k_1,
 \tag{4.10}
\]
and hence $W_{13}=B_0-2R_{11}-R_{12}-R_{13}$.
The cases $k_1=2$ and $k_1=1$ give
$B_0\le-k_3$ and $B_0\le1-k_2-k_3$, respectively.
The latter is positive only at $(1,0,0,2)$, of degree 3, where $R_{12}=1$.
Three fourth-band factors with a first-band factor exceed $N$.
For $k_1\ge3$, the lower slab leaves only $(3,0,0,1)$, where $B_0=0$.
It suffices that
\[
 \varepsilon_A N^{\,2/3-3\beta}>1:
 \tag{4.11}
\]
Three first-band factors and a fourth below $1/3$ have exponent
less than $3\beta+1/3<1$.
Four first-band factors have a smaller upper bound.

If $k_1=0,k_2\ge2$, then $B_0\le0$.
If $k_1=0,k_2=1$, then $B_0=2-k_3$.
At degree 3, take the second-band prime as $p_1$ and order the remaining primes as $p_2\le p_3$.
Then $p_1p_2^2\le e<N$, and
\[
 p_1p_2p_3^2=ep_3\ge\varepsilon_A N^{1+\lambda}\ge N
              \quad\text{if }\varepsilon_A N^\lambda\ge1.
 \tag{4.12}
\]
This is a $D_{13}$ representation.
When $k_3=0$, there is also one $D_{12}$ representation, and these cover $k_3=0,1,2$.
At degree 4, $k_3\ge2$, since $\beta+\lambda+2/3>1$.
Finally suppose $k_1=k_2=0$.
Degree 3 gives $k_3\ge1$ and one $D_{11}$ representation.
For $k_3=1$, there is also one $D_{12}$ representation, and they cancel $4-k_3$.
At degree 4, all primes lie in the third band, where $3\lambda+1/3>1$ and $B_0=0$.
\end{proof}

Conditions (4.11)--(4.12) and $\varepsilon_A N>N^{2/3}$ hold eventually,
and the finite radial range is $1+\log\varepsilon_A/L<\sum t_i<1$.
Section~5 gives (5.31)--(5.32) and positive excess at degrees 5 through 11.
Extra $D_{13}$ representations lower $W_{13}$.

Define the weighted sequences
\[
 \mathcal B_j=\{N-e:e\in\mathcal E_j\},\qquad
 \mathcal E_j=\{e:R_j(e)>0,\ \varepsilon_A N<e<N\}.
 \tag{4.13}
\]
Each occurrence has weight $R_j(e)$.
Then
\[
S_j\le\sum_eR_j(e)\mathbf1_{N-e\ {\rm prime}}\le S(\mathcal B_j;\mathcal P(N),N^{1/2})+O(N^{1/2+o(1)}).
\tag{4.14}
\]
Removing source conditions, including $(m,Nd)=1$, enlarges these positive counts.
Their negative coefficients preserve the required lower bound.
The sieve counts primes $N-e\ge N^{1/2}$ with weight $R_j(e)$,
and by (4.6) smaller primes and endpoints cost $O(N^{1/2+o(1)})$.
A source coordinate $q\mid N$ may leave a composite $N-e=qm$, and this excess is $E_{\rm src,exc}$.
Restricting to $p_i\nmid N$ gives equality up to the same error,
so (4.14) has a vanishing D-unit error.

The lower slab is exact:
\[
 p<(1-\varepsilon_A)N
       \ \Longleftrightarrow\ e=N-p>\varepsilon_A N.
 \tag{4.15}
\]
Its complement contributes zero.
The source factors are eventually odd, so $e\ne N$, and the shifted-prime source has $e\le N-3$.
For the fixed target, set
\[
\begin{gathered}
 A=3,\quad C_b=334,\quad N_{\rm cut}=N(1-L^{-333}),\\
 \Sigma_G=(\varepsilon_A N,N_{\rm cut}],\quad
 A_G^{\rm cut}=A_G\mathbf1_{e\le N_{\rm cut}}.
\end{gathered}
 \tag{4.16}
\]
The cut is applied before Buchstab and AA37.
Restoring it once costs
\[
 R0:G1\le100(NL^{-333}+1),\qquad
 E^D_{\rm cut}\le{100(NL^{-333}+1)\over\mathcal W}
                  =O(L^{-331})+O(L^2/N).
 \tag{4.17}
\]
Discrepancies, membership and boundary faces use (4.16).
\subsection{The four sieve rows and their coefficients}

This subsection sets up the four sieve rows and their signed remainder coefficients.
Put $\mathcal B_G^{\rm cut}=2\mathcal B_{11}^{\rm cut}
+\mathcal B_{12}^{\rm cut}+\mathcal B_{13}^{\rm cut}$, in which the source weights occur once.
Three Buchstab steps give
\[
\begin{aligned}
 S(\mathcal B_G^{\rm cut},N^\mu)
 ={}&S(\mathcal B_G^{\rm cut},N^u)
 -\sum_{N^u\le a_1<N^\mu}S((\mathcal B_G^{\rm cut})_{a_1},N^u)\\
 &+\sum_{N^u\le a_2<a_1<N^\mu}
                   S((\mathcal B_G^{\rm cut})_{a_1a_2},N^u)\\
 &-\sum_{N^u\le a_3<a_2<a_1<N^\mu}
                   S((\mathcal B_G^{\rm cut})_{a_1a_2a_3},a_3).
\end{aligned}
 \tag{4.18}
\]
Here $a_i\in\mathcal P(N)$, and the cutoffs exclude their endpoints.
Also $S(\mathcal B_G^{\rm cut},N^{1/2})\le S(\mathcal B_G^{\rm cut},N^\mu)$.
The cells are endpoint-clipped parent intervals $[N^{t_i},N^{\rho t_i})$ with inherited lower labels.
For the full anatomy infimum $\Theta_c$, put
\[
 \ell_c=\Theta_c-\eta_L-\sum_{i=1}^rt_i,\qquad
 D_c=N^{\ell_c},\qquad u_c=D_c^{\sigma^2}.
 \tag{4.19}
\]
AA37 is applied at $D_c$, with relative tail labels.
The smooth divisor satisfies $d<D_c^\sigma$, $d\mid P(u_c;N)$.
All source forms use the following four rows.
\begingroup
\setlength{\tabcolsep}{3pt}
\begin{longtable}{@{}llllp{0.40\textwidth}@{}}
\toprule
Row & Prefix & AA37 & $2A_R$ & Base cutoff and retained cells \\
\midrule
G1-0 & $1$ & upper $F$ & $-1$ & $N^u$; all \\
G1-1 & $a_1$ & lower $f$ & $+1$ & $N^u$; all \\
G1-2 & $a_1>a_2$ & upper $F$ & $-1$ & $N^u$; all \\
G1-3 & $a_1>a_2>a_3$ & lower $f$ & $+1$ & $E=N^{t_3}$ after PR; $\Theta_c-\eta_L-t_1-t_2-(3+2\tau)t_3\ge0$ \\
\bottomrule
\end{longtable}
\endgroup
The first three residuals are at least
\[
 {19101\over32000}-\eta_L,\quad
 {19101\over32000}-\eta_L-\rho\mu,\quad
 {19101\over32000}-\eta_L-2\rho\mu.
 \tag{4.20}
\]
For G1-0--G1-2, the smallest completion coordinate satisfies
$x_3\le u$, so the active $w$ branch is at least $19101/32000$.
If its guard closes, or the parent crosses $\mu$, the $v_1(x_1)$
branch gives the same lower bound. Hence
$\Theta_c\ge19101/32000$.
Dividing by $u$ gives values above 30.
For the fourth row, put
$A_{\rm floor}=53/256-\eta_L-(3/2)\mu\tau$.
The largest anatomy coordinate is $x_1\in[t_1,\rho t_1]$, and the full anatomy contains
$v_1(x)=\min\left\{{1+x\over2},{32-34x\over39},
                         {57\over64}-{3x\over2}\right\},$
also across $\mu$.
The middle branch exceeds the first for $x\le25/107$ and the third for $x\ge175/1568$.
Since $175/1568<25/107$, it is redundant.
Hence
$\Theta_c\ge\min\left\{{1+t_1\over2},
                    {57\over64}-{3\over2}\rho t_1\right\}.$
The increasing branch is evaluated on the left and the decreasing branch on the right.
For $t_2\le t_1\le\mu$, both residual branches decrease.
Thus
\[
 \begin{aligned}
 \Theta_c-t_1-t_2
 &\ge\min\left\{{1\over2}-{3\over2}\mu,
                {57\over64}-\left(1+{3\over2}\rho\right)\mu-\mu\right\}\\
 &=\min\left\{{53\over256},
              {53\over256}-{3\over2}\mu\tau\right\}
 ={53\over256}-{3\over2}\mu\tau.
 \end{aligned}
\]
Subtracting $\eta_L+t_3$ and using the retained-cell condition gives $\ell_c\ge A_{\rm floor}-t_3$ and $\ell_c\ge2t_3$, whence
$\ell_c\ge{2A_{\rm floor}\over3}>{1379\over10000}>\ell_0,
 \qquad \ell_0={53\over412}-{1\over20000}.$
The exact condition is $\ell_c\ge2\rho t_3$.
It gives the lower-sieve range on the whole cell, including equality.
The linked level and orbit data are:
\begin{center}
\begin{tabular}{@{}lllp{0.24\textwidth}@{}}
\toprule
Row & Lower floor for $\ell_c$ & Orbit $M_c$ & Ratio and outside error \\
\midrule
G1-0 & $19101/32000-\eta_L$ & $1$ & $s>30$; AA37 \\
G1-1 & $19101/32000-\eta_L-\rho\mu$ & $1$ & $s>30$; AA37 \\
G1-2 & $19101/32000-\eta_L-2\rho\mu$ & $1$ or $2$ & $s>30$; AA37 \\
G1-3 & $\max(A_{\rm floor}-t_3,2t_3)$ & $1$, $2$ or $6$ & $s\ge2$; AA37 and (4.57) \\
\bottomrule
\end{tabular}
\end{center}
All rows use all four source forms and level (4.19),
with exact signed remainder coefficient (4.26).
The common remaining errors are in the final table.


For the cut source, define
\[
\begin{aligned}
 C(q;N)&=\sum_eA_G^{\rm cut}(e)\mathbf1_{q\mid N-e},&
 X_G&=\sum_eA_G^{\rm cut}(e),\\
 Y_G(q)&=\sum_{(e,q)=1}A_G^{\rm cut}(e),&
 \Delta_0(q)&=C(q;N)-X_G/\varphi(q),\\
 \Delta_P(q)&=C(q;N)-Y_G(q)/\varphi(q).&&
\end{aligned}
 \tag{4.21}
\]
For allowed odd $p\nmid Nb$, take $\omega_b(p^k)=p/(p-1)$ for $k\ge1$,
so that $\omega_b(p^k)/p^k=1/\varphi(p^k)$.
Prefix primes lie above the cutoff and original tail primes lie below it.
Thus $r_{\rm std}((\mathcal B_G)_b,c')=\Delta_0(bc')$, including $c'=1$ and repeated tail primes.
Mertens and $\sum_p(p-1)^{-2}<\infty$ verify the density hypotheses uniformly.

Fix a cell and tail labels, and let $O(q),C_c(q)$ be the ordered masked and unrestricted labelled coefficients.
Put $M_c=\prod_h|G_h|!$.
With the actual masks $K_T,K_P$,
\[
 \Xi_U(q)={1\over M_c}\sum\lambda_d(1-K_TK_P)
                  \mathbf1_{q=d\prod p_j},\qquad
 O(q)={C_c(q)\over M_c}-\Xi_U(q).
 \tag{4.22}
\]
Equal-bin permutations preserve the clipped prime sets, products, labels and summands.
Unequal bins enforce order.
For G1-3, the $a_3$-cutoff difference is removed first, leaving base $E=N^{t_3}$.
The factor $1/M_c$ and the prefix-diagonal correction are kept.

Put $\Phi=M_c^{-1}\sum C_c\Delta_P$,
$NC_{\rm sf}=\sum_{q\ {\rm sf}}O(q)(\Delta_0-\Delta_P)(q)$,
and $\Xi_A(q)=O(q)\mathbf1_{q\ {\rm not\ sf}}$.
Adding and subtracting $\Delta_P$ gives
\[
\begin{gathered}
 R_{AA}=\Phi+NC_{\rm sf}-\sum_q\Xi_U(q)\Delta_P(q)
                   +\sum_q\Xi_A(q)(\Delta_0-\Delta_P)(q),\\
 \Delta_0-\Delta_P=-{X_G-Y_G(q)\over\varphi(q)}.
\end{gathered}
 \tag{4.23}
\]
Allowed repeats are kept with these masks.
Composite sources may retain the comparator difference.
Theorem~\ref{thm:repeated-primes} bounds both squareful terms.
The tested congruence is exactly
\[
 |(\mathcal B_G^{\rm cut})_q|
 =\sum_eA_G^{\rm cut}(e)\mathbf1_{q\mid N-e}
 =\sum_eA_G^{\rm cut}(e)\mathbf1_{e\equiv N\pmod q}.
 \tag{4.24}
\]
Every actual modulus satisfies $(q,2N)=1$ and
\[
 q<N^{e_G},\qquad e_G:=\frac{153}{256}(1+\tau+\sigma)<0.598.
\]
By (4.28), $X>cN$, where $c=\varepsilon_A2^{-24}>0$.
Eventually $N^{e_G}<(cN)^{3/5}<X^{3/5}$.
This verifies the full support condition of Theorem~\ref{thm:distribution-input};
see \cite[Theorem~3.2, (H3)]{Z1}.

Primes below $u_c$ determine $d$ uniquely, and rough positions are at least $u_c$.
Their multiplicities $e_p$ allow at most $m!/\prod_pe_p!$ assignments, and masks remove assignments.
Consequently
\[
 \ell_0={53\over412}-{1\over20000}>0,\qquad
 M_\sigma=\left\lceil{3\over5\ell_0\sigma^2}\right\rceil,\qquad
 |\gamma_q|\le M_\sigma! =: C_\sigma,
 \tag{4.25}
\]
which covers original, unrestricted and grouped coefficients.
Fixed row conventions are included in $C_\sigma$,
and $N_{\rm lab}(\sigma)$ enters aggregate bounds once.
For 1-bounded slot sequences we extract $C_{\rm slot}=\prod_s|J_s|!$.
The D-unit formulas are
\[
 {(2A_R)\nu_c\over M_c\mathcal W}\sum_qC_c(q)\Delta_B(q),
 \qquad
 {(2A_R)\nu_c C_{\rm slot}\over M_c\mathcal W}
                  \sum_q{C_c(q)\over C_{\rm slot}}\Delta_B(q).
 \tag{4.26}
\]
Here $\nu_c=+1$ for upper AA37 and $-1$ for lower AA37.
We form $C_c$ before this sign.
The slot factorial and source weights enter once, and $K_G$ is omitted from the remainder.

\subsection{Preparation of the product source}

This subsection splits the product source into boxes and checks the hypotheses of the distribution input.
Order-3 Heath--Brown \cite{ref6} is applied pointwise to every source coordinate.
Then
\[
 n(\mathbf v)=\prod_{i=1}^sv_i=e,\qquad s\le6k\le24.
 \tag{4.27}
\]
There are $27+81+27+27=162$ patterns per row, giving
648 row/form/pattern instances.\phantomsection\label{p2:g1-external-count}
All source indicators are kept.
Put $\delta=L^{-334}$.
We split each dyadic shell $[2^b,2^{b+1})$, $b\ge1$, into
$J=\lceil\log2/\log(1+\delta)\rceil$ intervals
$[2^b2^{j/J},2^b2^{(j+1)/J})$, $0\le j<J$.
Their common relative width is $\Delta=2^{1/J}-1$.
For $\delta\le1/300$, $\delta/3\le\Delta\le\delta$.
The point 1 is kept as a singleton.
A boundary box crosses a source, window, ordering, cap or cut face.
An interior box has constant source indicators.
Per variable there are $O(L)$ shells and $O(L^{334})$ intervals,
so there are $O(L^{335})$ choices per variable.
The 24-variable fine count is $O(L^{8040})$; two further shells give $K_0=8042$.
The three internal slot-shell sums are counted in $A_{\rm work}-A'=4$.

Choose a nonunit cell, and put
$T:=T_{\rm II}:=N_{II}:=V_{\rm sel},\qquad q:=1+\Delta,\qquad
 C_{\rm sel}=[T,qT),\qquad h_N:=\Delta T.$
Thus $\delta T/3\le h_N\le\delta T$.
The exact interval counts are
\[
 \begin{aligned}
 I_N&:=(T,qT]\subset(T,2T],\\
 H_N&:=\#(C_{\rm sel}\cap\mathbb Z)
      =\lceil qT\rceil-\lceil T\rceil,\\
 N_I&:=\#(I_N\cap\mathbb Z)=\lfloor qT\rfloor-\lfloor T\rfloor.
 \end{aligned}
\]
Both differ from $h_N$ by less than 1.
The ambient count is $N_I$, and roughness stays in its coefficient.
For an integer endpoint $w=2^{b+j/J}$, $w^J=2^{bJ+j}$ gives $J\mid j$, so only $2^b,2^{b+1}$ are integer endpoints.
The factor $\mathbf1_{P^-(n)\ge z_0(N)}$, with $z_0(N)>2$, kills these endpoints.
Hence
$a_{\rm sel}(n)\mathbf1_{C_{\rm sel}}(n)
 =a_{\rm sel}(n)\mathbf1_{I_N}(n)
 \qquad(n\in\mathbb Z).$

\begin{lemma}[Common product boxes]\label{lem:g1-f}
Let $B$ be a retained interior box.
For its lower endpoints $V_i$, put $X=\prod_iV_i$.
Then
\[
 n(\mathbf v)\in\Sigma_G\quad(\mathbf v\in B),\qquad
 \varepsilon_A2^{-24}N<X\le N.
 \tag{4.28}
\]
The largest eligible HB scale in
$[N^{60\sigma},N^{1/3+14\sigma}]$ is chosen, with the first index in a tie.
Write $N_{II}$ for this lower scale and $M$ for the complementary product.
Every box is Class II, and eventually
\[
 MN_{II}\asymp N,\quad
 X^{60\sigma}\le N_{II}\le X^{1/3+15\sigma},\quad
 N\le X^{1+\sigma/20}.
 \tag{4.29}
\]
On its doubled complementary majorant support,
\[
 N-mn\ge NL^{-334}\ge X(\log X)^{-334}>0.
 \tag{4.30}
\]
The thresholds (4.32)--(4.32a) below imply LEN on every box.
\end{lemma}
\begin{proof}
Buchstab, AA37, (4.23) and the G1-3 base $E$ preserve the measure.
Prime-power replacement and logarithm freezing preserve the indicators.
Equation (4.27) gives the support, and the bounds $V_i\le v_i<2V_i$ and $s\le24$ give the scale.

Choose a coordinate with exponent in $[\alpha,1/3]$.
One of its six or fewer factors is at least $N^{\alpha/6}$.
Its lower endpoint lies between half this size and $N^{1/3}$.
Since $\alpha/6>60\sigma$, an eligible choice exists, and the prescribed choice is constant on each box.
Fubini gives
\[
 \sum_{\mathbf v\in B}\prod_ia_i(v_i)H(\prod_iv_i)
       =\sum_{m,n}\beta(m)\alpha_{II}(n)H(mn).
 \tag{4.31}
\]
The selected coefficient is rough Möbius, rough 1, or normalized rough logarithm.
The complementary convolution has at most 23 factors and bound $\tau_{23}$.
Each logarithm stays on its original factor.
Endpoint subtraction is used in the interval and auxiliary-coprimality lemmas of Theorem~\ref{thm:distribution-input}.
Narrow intervals, the relative-norm alternative and both cross-copy masks are kept.
These inputs include more than scalar divisor moments.

Put $\omega=\log_N(N/X)$, and impose
\[
 L\ge\max\{2611/\sigma,300000\}
 \tag{4.32}
\]
\[
 \sigma L\ge2^{130}(\log L+1).
 \tag{4.32a}
\]
For fixed $\sigma>0$ the latter holds eventually, and uniformity as $\sigma\to0$ is unnecessary.
Since $T\ge N^{60\sigma}$,
$60\sigma L\ge60\cdot2^{130}(\log L+1)
              >334\log L+\log6,$
so $h_N\ge T/(3L^{334})>2$.
Thus $T/(6L^{334})\le N_I$ and $N_I\le2\delta T$.
Also
$T^{2^{-126}}\ge\exp(60\sigma L/2^{126})\ge12L^{334},$
since $60\sigma L/2^{126}\ge960(\log L+1)>334\log L+\log12$.
Consequently
$N_I\ge\frac{T}{6L^{334}}
       \ge2T^{1-2^{-126}}\ge(2T)^{1-2^{-126}}.$
With $Y=\lfloor T\rfloor$ and $H=N_I$, the integer intervals coincide.
Then $(2Y)^{1-2^{-126}}\le H\le Y$, using $T\ge2$ and $\delta\le1/300$.
This gives the short-interval input for $k=2^{126}$.
Threshold (4.32) gives
\[
 \omega<{\log(2^{24}/\varepsilon_A)\over L}\le{\sigma\over101},
 \quad {1/3+14\sigma\over1-\omega}\le1/3+15\sigma,\quad
 {2\omega\over1-\omega}<0.02\sigma.
 \tag{4.33}
\]
Also $(1-\omega)(1+\sigma/20)\ge1$, and these prove (4.29) and the packing scale conversion.
The frequency stays $a=N$.

Put $x=L^{-334}$; for $x\le1/300$, $(1+x)^{24}-1\le25x$.
Let $1\le r\le23$ count nonunit complementary cells, with endpoints $V_i,W_i=V_i(1+\Delta)$.
Put $M=\prod V_i$ and $M_+=\prod W_i$.
Unit factors stay singletons.
There is a nonunit complementary factor.
At least three source coordinates have product at least $N^\alpha$.
Only one HB factor was selected.
Thus $1+x/3\le M_+/M\le(1+x)^{23}\le1+25x$.
The ambient length is $M_I=M_+-M$, with $xM/3\le M_I\le25xM$.
The fixed majorant $\psi_I(m)=\phi_0((m-m_0)/M_I)$ in condition C3
of Theorem~\ref{thm:distribution-input} doubles this interval about its midpoint.
Its upper extension is $M_I/2\le12.5xM$.
Since $nM_+\le N_{\rm cut}$ on an interior box,
$mn\le N_{\rm cut}+12.5xN\le N_{\rm cut}+25xN$ on the doubled support.
Hence $N-mn\ge(L-25)NL^{-334}\ge NL^{-334}$ for $L\ge26$.
The function $t/(\log t)^{334}$ increases for $\log t>334$, and as $X\le N$ this proves (4.30).
Threshold (4.32) also gives $NL^{-334}>1$ and $\log X>334$.

Write $X=N^S$, $V_{\rm sel}=N^{e_{\rm sel}}$ and $M=X/V_{\rm sel}$.
The geometry gives
\[
 M_I\ge\tfrac13L^{-334}M,
 \qquad
 S/2-e_{\rm sel}\ge1/6-(15+1/202)\sigma>0
\]
where $S\ge1-\omega$ and $e_{\rm sel}\le1/3+14\sigma$.
Thus $M_I/X^{1/2}\ge\tfrac13L^{-334}N^{1/6-(15+1/202)\sigma}\to\infty$,
and in particular $M_I\ge X^{1/2}$ eventually.
Also $M_I/X^{1/2}\ge\tfrac13L^{-334}N^{0.15}>1$ eventually.
By (4.33),
$(1-\omega)^{-334}
 \le\exp(334\omega/(1-\omega))\le1+4\cdot10^{-10}.$
Since $\log X=(1-\omega)L$,
$L^{-334}\le(\log X)^{-334}\le(1+4\cdot10^{-10})L^{-334}$.
Therefore
\[
\begin{aligned}
 \tfrac17T(\log X)^{-334}&\le N_I\le2T(\log X)^{-334},\\
 \tfrac14M(\log X)^{-334}&\le M_I\le25M(\log X)^{-334}.
\end{aligned}
\]
As $M_+=Mq^r$ and $q^r\to1$ uniformly,
$m_0=M(1+q^r)/2\le2M=2X/T$ eventually.
On $[M,Mq^r]$, $\phi_0((m-m_0)/M_I)=1$, and both Cauchy--Schwarz copies remain in the doubled rectangle after coloring.
Equation (4.30) excludes exact resonance; see \cite[(3.10)]{Z1}.

By the internal Lemma~\ref{thm:packing-input},
\[
\begin{gathered}
 Q\le TN^{-(1+7271/256)\sigma}\le TX^{-(1+7271/256)\sigma},\\
 {7271\over256}-{1\over3}={21557\over768}=28+{53\over768}>28.
\end{gathered}
\]
Thus $QX^{\sigma/3}<TX^{-28\sigma}$.
Threshold (4.32a) gives $14\sigma L>334\log L+\log6$.
Since $\log X\ge L/2$,
$QX^{\sigma/3}<TX^{-28\sigma}
 \le T/(6L^{334})\le N_I.$
The same threshold gives $334\log L+\log3<\sigma L/100$.
Hence
$\log_NM_I\ge{2\over3}
 -\left(14+{1\over101}+{1\over100}\right)\sigma
 >{2\over3}-14.02\sigma.$
Packing also gives $\log_N(QRS)\le153/256+\sigma$.
Since $104.02\sigma<53/768=2/3-153/256$ and $X^{-89\sigma}\ge N^{-89\sigma}$,
$QRS<M_IX^{-89\sigma}$.
Finally, eventual monotonicity gives $\zeta=z_0(N)\ge z_0(X)$.
\end{proof}

Put $L_N=\log N$, $z=z_0(N)$.
The replacement $\mathbf1_{\mathbb P}(x)\mapsto\Lambda_z(x)/\log x$ costs (4.58)--(4.61).
Order-three rough Heath--Brown \cite{ref6} is used on each coordinate shell,
with the inverse Dirichlet series truncated at its upper endpoint.
The pointwise factors are 1 or $z$-rough.
The product inequalities are retained; they are constant on interiors.
Boundary pieces are estimated below.

For a form with $k\in\{3,4\}$ coordinates, let $H_B(\mathbf v)$
denote the HB numerator, including its source-form weight but excluding
the $k$ inverse coordinate logarithms. Then
$|H_B(\mathbf v)|\le2\cdot3^kL_N^k\le162L_N^4$.
Each of these $k$ HB logarithms is normalized by $L_N$, and the frozen coordinate logarithm satisfies $\lambda_{B,j}\ge\alpha L_N-1$.
For $L_N\ge2/\alpha$, the resulting exterior coefficient is at most
$2\cdot3^k(2/\alpha)^k\le162(2/\alpha)^4$ uniformly.
For a selected logarithm,
$\frac{\log n}{L_N}=\frac{\log X}{L_N}\frac{\log n}{\log X},
 \qquad \frac{\log X}{L_N}\le1.$
The selected coefficient is 1-bounded and the complement has bound $\tau_{23}$.
With (4.66) and fixed source-label/slot constants, this gives H7, $K_0=8042$.

The internal Lemma~\ref{thm:packing-input} is applied at $\Theta_c$.
The $d$ of AA37 is placed in slot 1, also when $J_1=\varnothing$.
Empty prime-label slots 2 and 3 carry unit atoms, and those two slots are rough.
The four surplus margins are
\[
 \left({7271\over256},52,{95221\over1280},{22407\over256}\right)\sigma;
 \tag{4.34}
\]
passing to $X\le N$ preserves the first margin.
After subtracting $1/3$ it is $28+53/768>28$, and the other margins exceed the losses in (4.33).
The box discrepancy is
\[
 \Delta_B(q;N)=\sum_{\mathbf v\in B}w(\mathbf v)
 \left[\mathbf1_{\prod v_i\equiv N\ (q)}
       -{\mathbf1_{(\prod v_i,q)=1}\over\varphi(q)}\right].
 \tag{4.35}
\]
For every source form and row, the preparation data are:
\begin{center}
\begin{tabular}{@{}lp{0.68\textwidth}@{}}
\toprule
Object & Common prepared value \\
\midrule
Selected factor & Largest eligible lower scale, first index in a tie; $T$ in (4.29) \\
Selected interval & $I_N=(T,qT]$, $N_I=\lfloor qT\rfloor-\lfloor T\rfloor$ \\
Complement & $\tau_{23}$; $M_I=M(q^r-1)$, $1\le r\le23$ \\
Majorant & $\phi_0((m-m_0)/M_I)$, $m_0\le2M$; doubled support (4.30) \\
External multiplicity & $648$ fixed instances; $O(L^{8042})$ boxes, (4.66) \\
\bottomrule
\end{tabular}
\end{center}
The hypothesis checks are collected below.
\begin{center}
\begin{tabular}{@{}ll@{}}
\toprule
Hypothesis & Verification \\
\midrule
H1, H2 & (4.22)--(4.26), with repeated-prime multiplicities \\
H3, LEN & Full support after (4.24), Lemma~\ref{lem:g1-f}, (4.34)--(4.35) \\
H4 & $a=N$, $K_{\rm scale}=31$, $(q,N)=1$ \\
H5 & The two windows of Lemma~\ref{lem:g1-f} \\
H6, H7 & (4.43), (4.63), (4.66) \\
\bottomrule
\end{tabular}
\end{center}
For H4, eventually $0<N\le X\log^{31}N\le X^{1+\sigma/20}$,
so Theorem~\ref{thm:distribution-input} applies to the plain cut product measure.
The local estimate is \cite[Theorem~3.2, (3.13)]{Z1}.
The spectral frequency is $a_T=\operatorname{sgn}(y)N$.
If $u_1,v_1$ invert $zb$ modulo $p_0c,c$, then $u_1-v_1\equiv0\pmod c$.
Hence $\operatorname{sgn}(y)N|y|(u_1-v_1)/c$ is integral, also for $p_0\mid c$, without any condition $(p_0,c)=1$.

\subsection{Source errors and the Class-II sum}

This subsection bounds the boundary, proper-power, logarithm and comparator errors and the interior Class-II sum.
A source face is
\[
g(\mathbf v)=\sum_ia_i\log v_i=\xi,\qquad
 |a_i|\le2,\quad\sum_i|a_i|\le48.
 \tag{4.36}
\]
Every defining-face coefficient has absolute value at most 2 after
expansion into the HB variables. Since $s\le24$, its coefficient sum
is at most 48.
The face families are:
\begin{center}
\begin{tabular}{@{}p{0.68\textwidth}l@{}}
\toprule
Family & Number \\
\midrule
Two $\Sigma_G$ faces, source windows/orderings, fixed coordinate transports & $\le16$ fixed \\
Coordinate dyadic edges & $\le4$ lattices \\
M\"obius caps $v_i=2X_m^{1/3}$ & $\le12$ lattices \\
Identity caps $n_m=(2X_m^{1/3})^3$ & $\le4$ lattices \\
Product dyadic edge and three slack families & $\le4$ lattices \\
\bottomrule
\end{tabular}
\end{center}
Thus there are at most 24 lattice families, of gap $\gamma=(\log2)/3$.
HB partition edges are refined exactly, and prefix and tail labels are modulus-side conditions.
The Class-II choice is constant on each box.

\begin{lemma}[Lattice thin shells]
Let $N\ge2$, $L=\log N$, $s\ge2$, and $a\in\mathbb Z^s\setminus\{0\}$.
Fix a nonempty $J_0\subseteq\{i:a_i\ne0\}$, and assume $|a_j|\in\{1,2\}$ for $j\in J_0$.
Let $1\le Y_0\le N$ and $0<\eta\le1/4$, and let $\Lambda\subset\mathbb R$ have pairwise distances at least $\gamma>0$.
Put $\kappa=(1-e^{-\gamma/2})^{-1}$.
Then
\[
\begin{aligned}
 &\#\{\mathbf v:\prod v_i\le N,\
 \operatorname{dist}(\sum a_i\log v_i,\Lambda)\le\eta,\
                         \max_{j\in J_0}v_j\ge Y_0\}\\
 &\quad\le s\left[4\eta\kappa N(1+L)^{s-1}
       +(2L/\gamma+2){N\over Y_0}(1+L)^{s-2}\right].
\end{aligned}
 \tag{4.37}
\]
For empty $\Lambda$, the count is zero.
For a singleton one may take $\kappa=1$ and interval-count factor 1.
\end{lemma}
\begin{proof}
Assign each tuple to the least $j\in J_0$ with $v_j\ge Y_0$, and apply \cite[Lemma~B.2, (B.14)]{Z1} on that fibre.
Only $|a_j|\in\{1,2\}$ matters there; the other coefficients enter $c=\sum_{i\ne j}a_i\log v_i$.
Put $Z=N/P$, where $P$ is the complementary product.
After clipping to $[Y_0,Z]$ and merging touching or overlapping intervals, the union has at most $2L/\gamma+2$ components.
Its fibre count is at most $4\eta\kappa Z+2L/\gamma+2$, the count before relaxation in (B.14).
Summing uses
\[
 \sum_{n\le y}\tau_k(n)\le y(1+\log y)^{k-1},\qquad
 \sum_{n\le y}{\tau_k(n)\over n}\le(1+\log y)^k,
 \tag{4.38}
\]
with $k-1$ factors fixed and the product constraint dropped in the harmonic sum.
These give (4.37).
The component count concerns the clipped union, and no restriction relates $\eta$ and $\gamma$.
\end{proof}

The boundary distance is at most $\eta=48L^{-334}$.
For each defining face, take $J_0=\{i:a_i\ne0\}$.
Six or fewer factors of a coordinate at least $N^\alpha$ give $Y_0=N^{\alpha/6}$.
A M\"obius-cap variable is at least $N^{\alpha/3}e^{-\eta}\ge Y_0$.
The same 16 fixed and 24 lattice families suffice.
The worst lattice has $\kappa=(1-2^{-1/6})^{-1}<9.2$, and $2L/\gamma+2\le10L$ for $L\ge2$.
Thus the boundary union $U$ for one form and pattern satisfies
\[
\begin{aligned}
 \#U\le{}&1091174.4\,NL^{-334}(1+L)^{23}\\
 &+24(16+240L)N^{1-\alpha/6}(1+L)^{22}\\
 \le{}&1091175\,NL^{-334}(1+L)^{23}\qquad(L\ge300000).
\end{aligned}
 \tag{4.39}
\]
The leading constant is $24\cdot192(16+24\cdot9.2)=1091174.4$.
The second-line ratio is bounded by
\[
 R(L)=5760e^{-(100/6894)L}L^{334}<3/5\quad(L\ge300000).
 \tag{4.40}
\]
Its logarithmic derivative is negative there.
At 300000, $e>8/3$, $(8/3)^{13}>300000$ and $(8/3)^9>5760$ give $\log R<9-(100/6894)300000+334\cdot13<-3/5$.
Since $e^{3/5}>5/3$, (4.40) follows.
This justifies the rounding at width exponent 334.

Boundary estimates keep the unnormalized HB logarithms.
The full weight $w(\mathbf v)=H_B(\mathbf v)\prod_{j=1}^k(\log p_j)^{-1}$
uses the actual coordinate logarithms; since $\log p_j\ge\alpha L\ge1$,
the numerator bound above gives the deliberately coarser bound below.
Thus for one pattern, including its source-form weight,
\[
 |w(\mathbf v)|\le162L^4,\qquad
 \sum_{q\le N^{3/5}}|\gamma_q|
  \left(\mathbf1_{q\mid N-e}+{1\over\varphi(q)}\right)
 \le C_\sigma\{\tau(N-e)+2(1+L)\}.
 \tag{4.41}
\]
Here $N-e\ge1$.
The reciprocal-totient constant is below 2, since
$q/\varphi(q)=\sum_{d\mid q}\mu^2(d)/\varphi(d)$ and
$\prod_p(1+1/(p(p-1)))<2$.
Put $F_*(e)=\tau(N-e)+2(1+L)$.
The composition bound is $\tau_a\tau_b\le\tau_{ab}$:
one maps the two compositions to a fixed matrix with those margins.
Cauchy--Schwarz and (4.38) give
\[
\begin{aligned}
 \sum_{e<N}\tau_{24}(e)\tau(N-e)^2
 &\le\left(\sum\tau_{576}\right)^{1/2}
       \left(\sum\tau_{16}\right)^{1/2}
 \le N(1+L)^{295},\\
 \sum_{e<N}\tau_{24}(e)F_*(e)^2&\le3N(1+L)^{295}
                              \quad(N\text{ sufficiently large}).
\end{aligned}
 \tag{4.42}
\]
Another Cauchy--Schwarz application on $U$ gives
\[
\begin{aligned}
 |E_{\rm bdry}|&\le162\sqrt{3\cdot1091176}\,C_\sigma\\
 &\qquad\cdot N(1+L)^{163}L^{-167}\\
 &<293108\,C_\sigma\\
 &\qquad\cdot N(1+L)^{163}L^{-167}.
\end{aligned}
 \tag{4.43}
\]
We include 648 instances and the multiplier $1/u=500$, since $u=1/500$;
we also include $N_{\rm lab}(\sigma)$ and $\mathcal W^{-1}$.
Since $C(N)\ge0.66016$,
\[
\begin{aligned}
 E^D_{\rm HB:G1}&\le3.6\cdot10^{10}C_\sigma N_{\rm lab}(\sigma)
                     (1+L)^{163}L^{-165}\\
 &\le3.6\cdot10^{10}C_\sigma N_{\rm lab}(\sigma)
                         e^{163/L}L^{-2}=o_\sigma(1).
\end{aligned}
 \tag{4.44}
\]
This covers cut faces, windows/orderings, dyadic families and caps.
The four parts count row/form/pattern/box occurrences once, including box and dyadic factors.
Hence H6 has $c_R=163$, with D rate $O_\sigma(L^{165-C_b/2})$.
Every fixed $C_b>330$ suffices; $C_b=334$ gives $O_\sigma(L^{-2})$.

The proper-power event concerns one source coordinate,
while other coordinates may contribute further prime factors.
For $n\ge N^\alpha$,
\[
 \mathbf1_{\mathbb P}(n)-{\Lambda(n)\over\log n}
       =-{1\over a}\mathbf1_{\substack{n=p^a\\a\ge2}}.
 \tag{4.58}
\]
Roughness below $N^\alpha$ preserves primes and removes powers.
The same absolute bound is used, telescoping over four or fewer coordinates.
For $Y=N^\alpha$,
\[
 \sum_{\substack{Y\le p^a\le N\\a\ge2}}p^{-a}
 \le2Y^{-1/2}+2(\log_2N)Y^{-2/3}\le3Y^{-1/2}
 \tag{4.59}
\]
Squares are bounded by an integer tail, and for $a\ge3$ we use $O(Y^{-(a-1)/a})$.
Fixing one coordinate at $t$ leaves at most
\[
 {N\over t}(1+\log N)^{\max(r-2,0)}\qquad(1\le r\le4)
\]
completions.
For $r=1$ the sole completion is bounded by $N/t$.
For $r=2$ we count directly, and for $r\ge3$ we use (4.38).
Thus $T_{PP}\le12N^{1-\alpha/2}(1+L)^2$.
By (4.25), the full modulus sum is
$O_\sigma(\tau(N-e)+2(1+L))$ per tuple.
Since $1\le N-e<N$, we use $\tau(N-e)\ll_\delta N^\delta$.
Then
\[
 |E_{PP}|\ll_{\sigma,\delta}N^{1-\alpha/2+\delta}(1+L)^3.
 \tag{4.60}
\]
We take $\delta=\alpha/8$ and absorb logarithms in the spare power.
Dividing by $\mathcal W$ gives
\[
 E^D_{PP:G1}\ll_\sigma N^{-\alpha/4}L^5
      =N^{-25/1149}O_\sigma(L^5)=o_\sigma(1).
 \tag{4.61}
\]
This step precedes HB boxes.
The event $p^2\mid q$ has its separate modulus bound.

For logarithm freezing, each coordinate has at most six HB factors.
If $\lambda_{B,m}$ is its lower logarithm,
\[
 \left|{1\over\log p_m}-{1\over\lambda_{B,m}}\right|
 \le {6L^{-334}\over(\alpha L-1)^2}.
 \tag{4.62}
\]
Up to four inverse logarithms change by at most $c_\alpha L^{-336}$.
This is multiplied by the unnormalized numerator bound $162L^4$ above.
The frozen inverse logarithms contribute to (4.62),
while the normalized interior coefficient $162(2/\alpha)^4$ excludes them.
Summing exact differences over disjoint boxes,
with $\sum\tau_{24}\le N(1+L)^{23}$, (4.42) and Cauchy--Schwarz, gives
\[
 |E_{LOG}|\ll_\sigma N(1+L)^{161}L^{-334},\qquad
 E^D_{LOG}\ll_\sigma L^{-171}\ll_\sigma L^{166-C_b}=L^{-168}.
 \tag{4.63}
\]
Hence H6 holds with $C_{\rm box}=C_b=334$ and $c_{\log}=164$.

For the comparator error, source primes are at least $N^\alpha$ before HB.
Thus
\[
 X_G-Y_G(q)
 \le100\sum_{\substack{p\mid q\\p\ge N^\alpha}}{N\over p}
 \le100\omega(q)N^{1-\alpha}.
 \tag{4.64}
\]
Also $\sum_q|O(q)|\omega(q)/\varphi(q)\ll_\sigma L^3$, by (4.25) and reciprocal-totient summation.
Equations (4.23)--(4.24) give
\[
 E^D_{NC:G1}\ll_\sigma N^{-\alpha}L^5=o_\sigma(1).
 \tag{4.65}
\]
The two squareful terms in (4.23) contribute
$O_\sigma(N^{-\ell_0\sigma^2}L^4)$ once, by Theorem~\ref{thm:repeated-primes}.
Class II includes GR cross-copy restoration.

The external precisions are
\[
\begin{gathered}
 A_{\rm raw}=3,\quad K_0=24C_b+26=8042,\\
 A'=A_{\rm raw}+K_0+2=8047,\quad
 A_{\rm work}=A'+4=8051.
\end{gathered}
 \tag{4.66}
\]
These use the \hyperref[p2:g1-external-count]{box count above}.
Three internal slot-shell sums use three of the four extra logarithms.
Choose
\[
\begin{aligned}
 A_E&\ge16770+c_*,\qquad C_1=A_E+2c_*+3,\\
 C_M&=5C_1+A_E+334+c_*+3,\qquad A_3=A_E+336,\\
 A''&=A_{BFI}+(334+A_3+1)/2.
\end{aligned}
 \tag{4.67}
\]
Here $16770=2A_{\rm work}+2C_b$.
Fix $c_*\ge\max\{c_{\rm pre}+2,c_1,c_{\rm diag},c_1(B_0),0\}$
and above all remaining preceding logarithmic losses, with
$c_{\rm pre},c_1,c_{\rm diag}$ as in (1.29).
Then (4.67) implies every lower bound in (1.29), in particular
$A_E\ge16772+c_{\rm pre}$.
Precisions are fixed before the threshold.
Theorem~\ref{thm:distribution-input} gives
\[
 |S_B|\ll X(\log X)^{-8047}
      +X^{1-c_\sigma}(\log X)^C,\qquad
 c_\sigma={\ell_0\sigma^2\over4}>0.
 \tag{4.68}
\]
This includes local restoration and GR at precision $A_{\rm work}$.
The prime-source condition $C_{\rm box}\ge\lceil2A+32\rceil$ uses
the global target $A$ of \cite[Theorem~1.2]{Z1}.
The local target $A'$ does not enter this box-width requirement.
The product-source H6 exponents are those above.
Lemma~\ref{lem:external-classii-sum} applies with $C_{\mathscr F}=C$:
by (4.2) and (3.1), $\mathcal W=W_N$, and (4.28) gives $X\ge N/L^{31}$ eventually.
Set $C_{\rm pow}=C+31c_\sigma$, which is a fixed finite exponent.
Hence
\[
 E^D_{II:G1}\ll_\sigma L^{-3}
 +N^{-c_\sigma}L^{8042+C_{\rm pow}+2}=o_\sigma(1).
 \tag{4.69}
\]
Scale conversion preserves $c_\sigma$, with no $N^{-\sigma}$ saving.
The stronger (4.28) is not needed here.
Boundary, comparator, proper-power and logarithm sums have no $L^{8042}$ factor.
The raw boundary rate is $O_\sigma(NL^{-4})$, with D rate $O_\sigma(L^{-2})$.
All these errors vanish at the fixed target.

\subsection{The dynamic-cutoff peel}

This subsection bounds the error of the dynamic-cutoff peel.
For a retained G1-3 cell, put $b=a_1a_2a_3$ and $E=N^{t_3}$.
The finite union bound is
\[
 S((\mathcal B_G)_b,a_3)\ge S((\mathcal B_G)_b,E)
       -\sum_{\substack{E\le r<a_3\\r\ {\rm prime}}}
                             S((\mathcal B_G)_{br},E).
 \tag{4.45}
\]
Here $a_1>a_2>a_3$, $r<a_3$, $r\nmid N$ and $(br,N)=1$.
The last two follow from $\mathcal P(N)$ and the coordinate conditions.
The negative upper-chain sign permits enlargement of this nonnegative difference.
Only the $E$-base enters HB.

\begin{lemma}[Uniform peel bound]
Let $0<\tau<8\cdot10^{-6}$ and $\rho=1+\tau$.
Suppose $(m,N)=1$ and $P^-(m)\ge z$, with
\[
 \log_Nm\le153/256+2\tau,\quad
 N^{u/\rho^3}\le z\le N^{\rho\mu},\quad w=\min(z,N^\alpha).
 \tag{4.46}
\]
Eventually, uniformly in $m$,
\[
 S((\mathcal B_G)_m,z)\le3\cdot502002\,{U_0\over m},
 \qquad U_0=2U_C.
 \tag{4.47}
\]
\end{lemma}
\begin{proof}
Write $N-e=mk$.
By the representation bound above, $A_G^{\rm cut}\le3$.
Put $\mathcal E_h=\{e:A_G^{\rm cut}(e)\ge h\}$ for $h=1,2,3$.
Then $A_G^{\rm cut}=\sum_{h=1}^3\mathbf1_{\mathcal E_h}$.
Each layer has the same roughness and two-form bound.
The layers are summed once, including coefficient 2 on both $R_{11}$ forms.

Enlarge the range to $1\le k\le\lfloor(N-1)/m\rfloor$, of length $K\le N/m$;
both $k$ and $N-mk$ are $w$-rough.
If $p<w$, $p\mid N$ and $p\mid k$, then $p\mid e$, which source roughness excludes.
Thus the all-primes two-form majorant applies.
For squarefree $d\mid P(w)$, CRT gives
\[
\begin{gathered}
 \nu_N(p)=1\ (p\mid N),\quad \nu_N(p)=2\ (p\nmid N),
 \quad \nu_N(2)=1,\\
 A_d=K{\nu_N(d)\over d}+R_d,\quad |R_d|\le\nu_N(d).
\end{gathered}
 \tag{4.48}
\]
For $p<w$, $p\nmid m$ follows from $P^-(m)\ge z\ge w$.
For $p\ge w$, the same rules are extended only in the Euler product.
A prime endpoint costs $g(w)=O(1/w)$ in the auxiliary sum.
Selberg least common multiples stay on $P(w)$.
Put $f(p)=\nu_N(p)/p$ and $g(p)=f(p)/(1-f(p))$.
They are extended squarefree multiplicatively, and we put
$G_N(w)=\sum_{d\le w,d\mid P(w)}g(d)$.
Finite-poset diagonalization gives
\[
 \sum_{d,e}\lambda_d\lambda_ef([d,e])
 =\sum_r{y_r^2\over g(r)},\quad
 y_r=\sum_{r\mid d}\lambda_df(d),\quad
 1=\lambda_1=\sum_r\mu(r)y_r.
 \tag{4.49}
\]
Cauchy--Schwarz gives minimum $1/G_N(w)$ at $y_r=\mu(r)g(r)/G_N(w)$.
Finite inversion gives
\[
 \lambda_d=\mu(d)\prod_{p\mid d}(1-f(p))^{-1}
     {\sum_{r\le w/d,(r,d)=1}g(r)\over G_N(w)},\\
 |\lambda_d|\le\prod_{p\mid d}(1-f(p))^{-1}
              \ll(d/\varphi(d))^2.
 \tag{4.50}
\]
The Dirichlet series factors as $\zeta(1+s)^2H_N(s)$,
where $H_N(s)=\sum h_N(r)r^{-s}$.
Its nonconstant local coefficients are
$g(p)-2/p$, $p^{-2}-2g(p)/p$ and $g(p)/p^2$.
For $p\nmid2N$, absolute logarithmic moments through order 2 converge.
For $p\mid N$, the absolute local mass is $O(1/p)$.
We split at $Y=\log(3N)$.
Mertens gives $\sum_{p\mid N,p\le Y}1/p\ll\log\log Y+1$.
Above $Y$ we use $\sum_{p\mid N}\log p\le\log N$.
The $j$-th logarithmic moment is $O((\log Y)^{j-1})$, for $j=1,2$.
Expanding absolute local products gives a fixed $B$ with
\[
 \sum_r|h_N(r)|(1+\log r)^2\ll(\log\log(3N))^B,\qquad
 H_N(0)={1\over2C(N)}.
 \tag{4.51}
\]
At zero, 2 contributes $1/2$.
A generic odd prime contributes $(1-1/(p-1)^2)^{-1}$.
For $p\mid N$, this is multiplied by $(p-2)/(p-1)$.
Convolution with $\sum_{n\le y}\tau(n)/n=\frac12\log^2y+O(\log(2y))$
and the second logarithmic moment for the tail $r>w$ give, uniformly,
\[
 G_N(w)={\log^2w\over4C(N)}
       \left(1+O\left({(\log\log(3N))^B\over\log w}\right)\right).
 \tag{4.52}
\]
Thus the dimension-two convention is
$G_N(w)\prod_{p<w}(1-\nu_N(p)/p)=e^{-2\gamma}/2+o(1)$.
Equations (4.48), (4.50) give
\[
 S([1,K],w)\le {K\over G_N(w)}+O(w^2\log^5w);
 \tag{4.53}
\]
the remainder is at most
$(\sum_{d\le w}|\lambda_d|\nu_N(d))^2
\ll w^2\log^2w(\log\log(3w))^4$.
Uniformly in (4.46),
\[
 \log_Nw\ge2/1001,\qquad
 w^2m\le N^{253/256+(153/64)\tau}<N^{1-0.0116}.
 \tag{4.54}
\]
These use $(1+8\cdot10^{-6})^3<1001/1000$, $u/\rho^3\ge2/1001$ and $\alpha>2/1001$,
together with $w\le z\le N^{\rho\mu}$ for the uniform upper saving.
For one layer, the main coefficient in $U_0/m$ units is eventually at most
$2.001/(2/1001)^2=501251.00025<502002$.
The strict slack absorbs the remainder in (4.53).
This proves (4.47) on the $N$-scale, including short $K$.
\end{proof}

For $m=br$ in (4.45), the retained-cell condition gives
\[
 \begin{aligned}
 \log_Nm&\le\rho(t_1+t_2+2t_3)\\
 &\le\rho[\Theta_c-\eta_L-(2\rho-1)t_3]\\
 &\le\rho(153/256-\eta_L)<153/256+2\tau.
 \end{aligned}
 \tag{4.55}
\]
On G1-1, G1-2 and G1-3, $\Theta_c\le\overline H=153/256$.
The two-branch peak is $(1+\mu)/2=153/256$.
Masked branches are zero or below it, and extra $w$-branches are zero or at most $(153-96c)/256\le153/256$.
The inherited label has $N^{t_3}\le a_3<N^{\rho t_3}$,
even for a clipped parent interval.
Since $N^u\le a_3<N^\mu$,
$t_3>u/\rho\ge u/\rho^3$ and $t_3<\mu\le\rho\mu$.
Hence $E\in[N^{u/\rho^3},N^{\rho\mu}]$.
All prime factors of $m$ are at least $E$, and $(m,N)=1$.
We apply (4.47) with $z=E$; $w$ may be smaller.
For fixed positive $\tau$, Mertens gives eventually
\[
 M_h=\sum_{p\ {\rm in\ parent\ bin}\ h}{1\over p}\le2\tau,
 \qquad\sum_hM_h\le4.59,
 \tag{4.56}
\]
Here $\log(\mu\rho/u)<4.59$.
The pair $a_3,r$ belongs to one parent bin. Thus
\[
\sum(a_1a_2a_3r)^{-1}\le4.59^2\sum_hM_h^2
\le2\cdot4.59^3\tau.
\]
The raw difference is below $3\cdot97100000\tau U_0$.
Its adverse D bound is
\[
 E^D_{\rm PR:peel}\le {3\cdot97100000\over2(1-\varepsilon_A)}\tau
 ={485500000000\over3333}\tau<145664566.4567\tau.
 \tag{4.57}
\]
The source multiplier enters once, without 500, label counts or tail factorials.
The slack includes the Selberg remainder.

\subsection{Canonical profiles and the joint lower bound}

This subsection compares the rows with canonical profiles and proves the joint lower bound for the four rows.
Put
$v(s)=v_1(s)-u,\qquad \psi(a,b)=a\mathbf1_{a\ge b},$
with the \hyperlink{p2:def:six-branch-saved}{six-branch definition}
of Section~\ref{p2:sec:2.1}, where $v_1$ is the unbuffered function of \eqref{p2:eq:row-level-function}.
Thus
\[
 \Theta_*(t_1,t_2,t_3)=\Theta_3^{\rm saved}(t_1,t_2,t_3).
 \tag{4.70}
\]
For $u\le t_3\le t_2\le t_1\le\mu$, put
\[
\begin{aligned}
v_1(s)&=\min\left\{\frac{1+s}{2},\frac{32-34s}{39},
                     \frac{57}{64}-\frac32s\right\},\\
v(s)&=v_1(s)-u,\qquad \psi(x,y)=x\mathbf1_{x\ge y}.
\end{aligned}
\]
The six branches are
\[
\begin{aligned}
\Theta_*(a,b,c)=\max\{&v(a),v(b),v(a+b),v(a+b+c),\\
 &\psi(v(a+c),a+2b+c),\psi(v(b+c),2a+b+c)\}.
\end{aligned}
\]
The $u$ subtraction precedes each $\psi$ guard.
For the actual coordinates, set
\[
 a_3={\Theta_*-t_1-t_2-t_3\over t_3},\qquad
 J=\int_u^\mu{dt_1\over t_1}\int_u^{t_1}{dt_2\over t_2}
                  \int_u^{t_2}{1-f(a_3)\over t_3^2}\,dt_3.
 \tag{4.71}
\]
Here $f$ is extended by zero below 2.
Above every finite table endpoint its true values are used.

The function $v_1$ is $3/2$-Lipschitz, and each exponent changes by at most $\mu\tau$ in a parent cell.
An active $\psi$-branch has unshifted gap $u-(9/2)\mu\tau-4\mu\tau>0$ throughout.
After $\eta_L$, its source margin is
\[
 g=(u-\eta_L)-(9/2)\mu\tau>0.001899>2\mu\tau.
 \tag{4.72}
\]
The other branches and maxima preserve this comparison across switches, and so do the two full-anatomy branches.
If a parent crosses $\mu$, then $t_1\le\mu\le\rho t_1$, so $|x_1-\mu|\le\mu\tau$ for its largest actual exponent.
Since $v_1(\mu)=153/256$,
$v_1(x_1)-\eta_L\ge153/256-(3/2)\mu\tau-\eta_L$.
Every canonical branch is at most $153/256-u$, so the margin is at least $u-\eta_L-(3/2)\mu\tau\ge g>0$.
Parents below $\mu$ use the activation argument above, and clipping shrinks the infimum domain.

If $f(a_3)>0$, the numerator exceeds $2t_3$,
and (4.72), for actual exponents then inherited labels, gives $\ell_c\ge2\rho t_3$.
Thus positive canonical lower terms lie in retained whole cells, and elsewhere their contribution is zero.
On retained cells we use Mertens and the monotonicity of $sf(s)$, keeping the numerator fixed.
The $E$-base dominates the actual-cutoff lower main term, so canonical restriction loses zero.
The arithmetic difference (4.57) remains, with possibly nonzero mass.

Let $H_{\log}=\log(\mu/u)$.
By (4.20), the first three residuals satisfy
\[
 0\le R\le500(F(30)-1)+500H_{\log}(1-f(30))
                         +250H_{\log}^2(F(30)-1)=:R_{\rm up}.
 \tag{4.73}
\]
Three elementary integrations give
\[
 500(1-H_{\log}+H_{\log}^2/2)
 -\int_u^\mu{dt_1\over t_1}\int_u^{t_1}{dt_2\over t_2}
                            \int_u^{t_2}{dt_3\over t_3^2}
 ={1\over\mu}
 \tag{4.74}
\]
which cancels the constant-one parts of $+F,-f,+F,-f$.
The actual upper main chain is therefore at most
\[
 G_1={1\over\mu}+R+J,\qquad0\le R\le R_{\rm up},
 \tag{4.75}
\]
with the preceding errors.

Let $X_j=\sum_eR_j(e)\mathbf1_{\varepsilon_A N<e<N}$,
and $X_j^{\rm cut}=\sum_eR_j(e)\mathbf1_{\varepsilon_A N<e\le N_{\rm cut}}$.
Then $X_G=2X_{11}^{\rm cut}+X_{12}^{\rm cut}+X_{13}^{\rm cut}$, and (4.17) bounds the mass difference.
PNT and partial summation give
\[
 X_j=\big((1-\varepsilon_A)I_j+o(1)\big){N\over L},
 \tag{4.76}
\]
where
\[
\begin{aligned}
 I_{11}={}&\int_\lambda^{1/3}
       {\log((1-2t)/t)\over t(1-t)}\,dt\\
 &+\int_\lambda^{1/4}{dt\over t}
       \int_t^{(1-t)/3}{\log((1-t-2v)/v)\over v(1-t-v)}\,dv,\\
 I_{12}={}&\int_\alpha^{1/3}{\log(2-3t)\over t(1-t)}\,dt,\\
 I_{13}={}&\int_\beta^\lambda
       {\log((1-t-\lambda)/\lambda)\over t(1-t)}\,dt.
\end{aligned}
 \tag{4.77}
\]
The last-prime density on the radial face is $1/(1-t-v)$.
Integrating the companion variable gives $I_{12},I_{13}$.
For $I_{11}$ the Buchstab argument is at most $214/93$, and its value is $v/(1-t-v)$ plus
$v\log((1-t-2v)/v)/(1-t-v)$ for $v\le(1-t)/3$.
This gives both integrals and the positive correction.
Set $e=yN$, $\varepsilon_A<y<1$.
On compact exponent boxes, the last logarithm has radial-face error
$O_{\varepsilon_A}(1/L)$, and $\int_{\varepsilon_A}^1dy=1-\varepsilon_A$.
PNT is applied first and fixed partitions are then refined,
and the positive exponent lower bound gives domination.
Monomial-face strips, including $D_{13}$'s endpoint-swap strip, have vanishing main mass.
Representation conditions are retained before the limit.
The upper cut costs (4.17).
This proves (4.76) at finite $\varepsilon_A$.

With $V_N(z)=(1+o(1))2e^{-\gamma}C(N)/\log z$, (4.75)--(4.76) give
\[
 {2S_{11}+S_{12}+S_{13}\over4(1-\varepsilon_A)U_C}
 \le K_GG_1+E^D_{G1},\qquad
 K_G={2I_{11}+I_{12}+I_{13}\over2e^\gamma}.
 \tag{4.78}
\]
The factor $1-\varepsilon_A$ cancels once here.
It remains in (4.57) and in the finite target of Section~1.
At $a=9/5$, that target is $8\log(50009/39991)$, with zero-epsilon limit $8\log(5/4)$.

The directed bounds are
\[
\begin{gathered}
 R_{\rm up}<4.4\cdot10^{-6},\qquad G_1\le\overline G,\\
 \overline G=6.0884403625974206828946573628908904\ldots<6.088440363.
\end{gathered}
 \tag{4.79}
\]
Appendix~A.5 gives the delay-function, quadrature and mass enclosures.
Its endpoint $K_{\rm hi}$ is \eqref{eq:g1-khi}.
Here $\overline G$ is the 4096-cell outward upper-corner sum in Appendix~A.5.
It is a conservative bound, not the continuous integral value.
Let $r_0,r_1,r_2,J_{\rm hi}$ denote the directed upper parts of $\overline G$.
Put
\[
\begin{aligned}
 T_0&=500+r_0,&T_1&=-500H_{\log}+r_1,\\
 T_2&=250H_{\log}^2+r_2,&
 T_3&=-(500-500H_{\log}+250H_{\log}^2-1/\mu)+J_{\rm hi}.
\end{aligned}
 \tag{4.83}
\]
Exactly, $\sum T_r=1/\mu+r_0+r_1+r_2+J_{\rm hi}=\overline G$.
Their signs are $+,-,+,-$, so $-K_GT_r$ has signs $-,+,-,+$.
We therefore use upper $K_G$ for rows 0,2 and lower $K_G$ for rows 1,3.
The four directed row enclosures are in Appendix~A.5,
and the four terms are cancelled before rounding.
The joint reference value is
\[
\begin{aligned}
 L_{G1}&=
 -{243950428265194656254206921939947762868871361874409552095097
 \over100433627766186892221372630771322662657637687111424552206336}\\
 &=-2.42897158741611974640561365334839795714\ldots .
\end{aligned}
 \tag{4.84}
\]
Put
$\delta_{G1}=4.98\cdot10^{-54},\qquad
 E_{G1}=L_{G1}-\delta_{G1},$
and let $R_{G1}$ sum the exact four row lower endpoints.
Appendix~A.5 proves $L_{G1}-R_{G1}<\delta_{G1}$, so $E_{G1}$ is used.
The main base contains the reference value $L_{G1}$ once,
and the adverse sum contains $\delta_{G1}$ once.
The row enclosures include $K_G$.
The cap for $G_1-2$ is $4.088440363$, which gives the weaker value $-K_{\rm hi}(2+4.088440363)=-K_{\rm hi}6.088440363$,
with about $1.61\cdot10^{-10}$ D loss.
Equations (4.13)--(4.78) transfer these enclosures to the arithmetic rows.

Successive differences are taken in the order
\[
 \begin{gathered}
 \text{source restoration, source conditions, PR-M, AA37,}\\
 \text{squareful/NC, proper powers, LOG, boundary, interior Class II}.
\end{gathered}
\]
AA37 includes fixed-parameter and asymptotic terms.
Only the $E$-base enters G1-3 Heath--Brown \cite{ref6}.
After monotonicity the exact first-contribution partition is retained,
and the favourable fixed-cutoff and G1 cutoff differences are discarded.
Each difference or part occurs once below.
\begingroup
\setlength{\tabcolsep}{3pt}
\begin{longtable}{@{}>{\raggedright\arraybackslash}p{0.24\textwidth}>{\raggedright\arraybackslash}p{0.68\textwidth}@{}}
\toprule
Contribution & Adverse D bound at fixed $\sigma,\tau$ \\
\midrule
Upper source restoration & (4.17) \\
Lower slab & zero: outside the source by (4.15) \\
G1 source-tag exceptions & $O_\sigma(N^{-\alpha/2+o(1)})$; the DIRECT $O_\sigma(L^3/N)$ term is separate, as in (1.23) \\
G1-3 peel & $145664566.4567\tau$, once \\
Matched canonical restriction loss & zero, by (4.72) \\
AA37 & its common $C_E\sigma$ share once, plus $o_\sigma(1)$ \\
Squareful terms & Theorem~\ref{thm:repeated-primes}: $O_\sigma(N^{-\ell_0\sigma^2}L^4)$ \\
NC & (4.65) \\
Tested proper powers & (4.61) \\
LOG & (4.63) \\
HB boundary & (4.44), one term with four parts \\
Interior Class II, including GR & (4.69) \\
Prime mass/PNT and Mertens & $o_\sigma(1)$, from (4.76)--(4.78) \\
\bottomrule
\end{longtable}
\endgroup
For $p\mid N$, there are $O((N/p)L^C)$ complementary products.
There are $O(L)$ such source primes, all at least $N^\alpha$.
A positive fraction of $\alpha$ absorbs divisor/logarithmic losses.
Labels and HB boxes preserve these sums.

Theorem~\ref{thm:distribution-input} and Theorem~\ref{thm:repeated-primes} in Appendix~B.1,
together with the directed endpoints above, give
\[
 \sum_{r=0}^3 Q_{G1-r,0}(N)
 \ge L_{G1}-\delta_{G1}-145664566.4567\tau-E^D_{AA:G1}(\sigma)
                    -E^D_{{\rm remaining}:G1}(N;\sigma),
 \tag{4.85}
\]
where the last term sums the vanishing errors in the table.
The AA37 term $C_E\sigma$ and the PR term $O(\tau)$ stay fixed.
Analytic constants and the residual margin are fixed first.
Then $\sigma$ is chosen, followed by $\tau=\sigma^9$, labels, widths and precisions.
Finally $N_0$ is chosen for all length and spectral conditions.
These include BFI/SW, Poisson and the W-prime bound at width 334.
The threshold depends on fixed $\sigma>0$.
Thus (4.85) holds for all four arithmetic rows, with $E_{G1}=L_{G1}-\delta_{G1}$.
